\documentclass[10pt,a4paper]{amsart}
\usepackage[margin=19mm]{geometry}
\usepackage{amsmath,amssymb,mathtools}
\usepackage[scr=rsfs,cal=euler]{mathalfa}
\usepackage{xfrac}
\usepackage[unicode,hidelinks,bookmarks=false]{hyperref}
\newcommand{\ie}{i.e.\ }
\newcommand{\cf}{cf.\ }
\newcommand{\resp}{respectively\ }
\newcommand{\Subsection}{\subsection}
\providecommand{\nobreakdash}{\nobreak\hbox{-}}
\setlength{\emergencystretch}{2em}
\multlinegap0pt
\usepackage{bm}
\usepackage{bbm}
\usepackage{dsfont}
\usepackage{xspace}
\usepackage{cancel}
\usepackage{mathtools}
\usepackage{amsthm}
\makeatletter
\@ifundefined{theorem}{\newtheorem{theorem}{Theorem}}{}
\makeatother
\title[Strong chain control sets and affine control systems]{Strong chain control sets and affine control systems}
\author{Fritz Colonius, Alexandre J. Santana}
\date{Source: arXiv:2407.18765v1 (CC BY 4.0)}
\begin{document}
\maketitle

\noindent\emph{Source and licence.} Strong chain control sets and affine control systems. Version arXiv:2407.18765v1 (CC BY 4.0), distributed under the Creative Commons Attribution 4.0 International licence, as recorded in the arXiv licence metadata for this exact version and shown on the abstract page https://arxiv.org/abs/2407.18765v1. Full rights notice: licenses/arxiv-2407.18765-CC-BY-4.0.txt.\par\smallskip

\noindent\textit{Editorial scope.} Counted unit: the Theorem whose source label is \texttt{Theorem\_equivalence} in the article (source0.tex lines 1040--1061, source label \texttt{Theorem\_equivalence}), delivered together with the article's own complete original proof of it (source0.tex lines 1063--1232, one \texttt{proof} environment of 170 lines). No mathematics is written, repaired or re-derived: statement and proof are the article's own text line for line. Required context retained uncounted: control\_affine (lines 524--527); metric\_U (lines 606--611). Every definition, display and label of those lines is retained unchanged. Omitted from this package and not redistributed: 1--523, 528--605, 612--1039, 1062--1062, 1233--2709. Nothing inside a retained excerpt is omitted or altered: every retained body is delivered exactly as the article's lines are written, so no omission receipt of any kind accompanies this package. Citations: the retained excerpts carry no citation command at all, so no bibliography entry of the article is transcribed and no reference is redistributed. Reference adaptation: none was needed and none was made; every reference inside the delivered unit resolves inside it, so no unresolved reference or citation is produced. Printed slips kept literal: no printed word, symbol, hypothesis or number of the counted statement or of its proof was altered. Layout and class: the article's own class is \texttt{elsarticle} with packages including amssymb, amsmath, amsthm, amsfonts, graphicx, enumerate, hyperref; the wrapper is the lane's pinned \texttt{amsart} editorial wrapper at 10pt on A4, which restates the theorem-like environments the retained text needs and adds the article's own macro definitions that the retained text needs (none), with extra wrapper packages (bm, bbm, dsfont, xspace, cancel, mathtools). The delivered numbering is the unit's own. Assets: no retained span contains a figure environment, a graphics-inclusion command, a PDF-page inclusion, a table or a TikZ picture; no image file is copied into this package, the package declares no asset dependency and no image file is redistributed with it. Licence: version arXiv:2407.18765v1 of this article is distributed under the Creative Commons Attribution 4.0 International licence, as recorded in the arXiv licence metadata for this exact version and shown on the abstract page https://arxiv.org/abs/2407.18765v1. A full rights notice is delivered at licenses/arxiv-2407.18765-CC-BY-4.0.txt. Characters: every retained excerpt is pure ASCII, so the delivered engine, XeLaTeX with its default Latin Modern text font, reports no missing character.
\begin{equation}
\dot{x}(t)=X_{0}(x(t))+\sum_{i=1}^{m}u_{i}(t)X_{i}(x(t)),\,u\in\mathcal{U},
\label{control_affine}%
\end{equation}

\begin{equation}
d_{\mathcal{U}}(u,v)=\sum_{i=1}^{\infty}\frac{1}{2^{i}}\frac{\left\vert
\int_{\mathbb{R}}\left\langle u(t)-v(t),y_{i}(t)\right\rangle dt\right\vert
}{1+\left\vert \int_{\mathbb{R}}\left\langle u(t)-v(t),y_{i}(t)\right\rangle
dt\right\vert }, \label{metric_U}%
\end{equation}

\begin{theorem}
\label{Theorem_equivalence}Consider a control system of the form
(\ref{control_affine}) on $M$ with control flow $\Phi$ on $\mathcal{U}\times
M$.

(i) If $E^{\ast}\subset M$ is a strong chain control set, then%
\begin{equation}
\mathcal{E}^{\ast}:=\{(u,x)\in\mathcal{U}\times M\left\vert \varphi(t,x,u)\in
E^{\ast}\text{ for all }t\in\mathbb{R}\right.  \} \label{lift_E}%
\end{equation}
is a maximal invariant strongly chain transitive set for the control flow
$\Phi$.

(ii) Conversely, let $\mathcal{E}^{\ast}\subset\mathcal{U}\times M$ be a
maximal invariant strongly chain transitive set for the control flow $\Phi$.
Then
\begin{equation}
E^{\ast}:=\left\{  x\in M\left\vert \exists u\in\mathcal{U}:(u,x)\in
\mathcal{E}^{\ast}\right.  \right\}  \label{down_E}%
\end{equation}
is a strong chain control set.
\end{theorem}

\begin{proof}
(i) Let $E^{\ast}$ be a strong chain control set and consider
$(u,x),\;(v,y)\in\mathcal{E}^{\ast}$ defined by (\ref{lift_E}). Pick $T>0$ and
$\varepsilon\in\mathcal{P}(\mathcal{U}\times M)$.

\textbf{Claim 1:} Define $\varepsilon_{M}(x):=\min\left\{  \varepsilon
(u,x)\left\vert u\in\mathcal{U}\right.  \right\}  ,x\in M.$ Then
$\varepsilon_{M}\in\mathcal{P}(M)$.

By compactness of $\mathcal{U}$ it is clear that $\varepsilon_{M}(x)>0$ for
all $x\in M$ and it remains to show that $\varepsilon_{M}$ is continuous.
Consider a sequence $x_{i}\rightarrow x$. It follows that $\varepsilon
(u,x_{i})\geq\varepsilon_{M}(x_{i})$ for all $u\in\mathcal{U}$. Continuity of
$\varepsilon$ implies that $\varepsilon(u,x)\geq\lim\sup_{i\rightarrow\infty
}\varepsilon_{M}(x_{i})$ for all $u\in\mathcal{U}$ and hence $\varepsilon
_{M}(x)\geq\lim\sup_{i\rightarrow\infty}\varepsilon_{M}(x_{i})$.

By compactness of $\mathcal{U}$ there are $u_{i}\in\mathcal{U}$ with
$f(u_{i},x_{i})=\varepsilon_{M}(x_{i})$ and there is a subsequence $(u_{i_{k}%
})$ converging to some $u^{\ast}\in\mathcal{U}$. It follows that
$\varepsilon_{M}(x_{i_{k}})=\varepsilon(u_{i_{k}},x_{i_{k}})\rightarrow
\varepsilon(u^{\ast},x)\geq\varepsilon_{M}(x)$. Since this applies to any
subsequence of $(u_{i})$ it follows that $\varepsilon_{M}(x)\leq\lim
\inf_{i\rightarrow\infty}\varepsilon_{M}(x_{i})$ proving \textbf{Claim 1}.

Strong chain controllability from $\varphi(2T,x,u)\in E^{\ast}$ to
$\varphi(-T,y,v)\in E^{\ast}$ yields the existence of a controlled
$(\varepsilon_{M},T)$-chain given by $x_{0}=\varphi(2T,x,u),x_{1},\ldots
,x_{k}=\varphi(-T,y,v)$ in $M,\,u_{0},\ldots,u_{k-1}\in\mathcal{U}%
,\,T_{0},\ldots,T_{k-1}\geq T$ with
\begin{equation}
d_{M}(\varphi(T_{j},x_{j},u_{j}),x_{j+1})<\varepsilon_{M}(\varphi(T_{j}%
,x_{j},u_{j}))\text{\ for }j=0,\ldots,k-1. \label{r1}%
\end{equation}
We construct an $(\varepsilon,T)$-chain from $(u,x)$ to $(v,y)$ in the
following way. Define
\[%
\begin{array}
[c]{cccc}%
T_{-2}=T, & x_{-2}=x, & v_{-2}=u, & \\
T_{-1}=T, & x_{-1}=\varphi(T,x,u), & v_{-1}(t)= & \left\{
\begin{array}
[c]{cc}%
u(T_{-2}+t) & \text{for }t\leq T_{-1}\\
u_{0}(t-T_{-1}) & \text{for }t>T_{-1}%
\end{array}
\right.
\end{array}
\]
and let the times $T_{0},\ldots,T_{k-1}$ and the points $x_{0},\ldots,x_{k}$
be as given earlier. Furthermore, set
\[%
\begin{array}
[c]{ccc}%
T_{k}=T, & x_{k+1}=y, & v_{k+1}=v,
\end{array}
\]
and define for $j=0,\ldots,k-2$%
\begin{align*}
v_{j}(t)  &  =\left\{
\begin{array}
[c]{lll}%
v_{j-1}(T_{j-1}+t) & \text{for} & t\leq0\\
u_{j}(t) & \text{for} & 0<t\leq T_{j}\\
u_{j+1}(t-T_{j}) & \text{for} & t>T_{j},
\end{array}
\right. \\
v_{k-1}(t)  &  =\left\{
\begin{array}
[c]{lll}%
v_{k-2}(T_{k-2}+t) & \text{for} & t\leq0\\
u_{k-1}(t) & \text{for} & 0<t\leq T_{k-1}\\
v(t-T_{k-1}-T) & \text{for} & t>T_{k-1},
\end{array}
\right. \\
v_{k}(t)  &  =\left\{
\begin{array}
[c]{ll}%
v_{k-1}(T_{k-1}+t) & \text{for }t\leq0\\
v(t-T) & \text{for }t>0.
\end{array}
\right.
\end{align*}
Define a constant $\widehat{\varepsilon}$ by%
\[
\widehat{\varepsilon}:=\frac{1}{2}\min_{j=-2,\ldots,k}\varepsilon_{M}%
(\varphi(T_{j},x_{j},v_{j})).
\]
Recall the definition of the metric $d_{\mathcal{U}}$ on $\mathcal{U}$ given
in (\ref{metric_U}) and choose $N\in\mathbb{N}$ large enough such that
$\sum_{i=N+1}^{\infty}2^{-i}<\widehat{\varepsilon}$. There exists $S>0$ such
that%
\[
\int_{\mathbb{R}\setminus\left[  -S,S\right]  }\left\vert y_{i}(\tau
)\right\vert \,\,d\tau<\widehat{\varepsilon}/\mathrm{\,diam}\,U\text{ for
}i=1,\ldots,N.
\]
We can assume without loss of generality that $T\geq S$. Since $\varepsilon
\in\mathcal{P}(\mathcal{U}\times M)$ and $T\geq S$ are arbitrary, the
following claim implies that $\mathcal{E}$ is strongly chain transitive.

\textbf{Claim 2: }An $(\varepsilon,T)$-chain from $(u,x)$ to $(v,y)$ is given
by%
\[
(v_{-2},x_{-2}),\,(v_{-1},x_{-1}),\ldots,(v_{k+1},x_{k+1})\text{ in
}\mathcal{U}\times M\text{ and }T_{-2},\,T_{-1},\ldots,T_{k}\geq T.
\]
For the proof of the claim observe first that there are no jumps at
$x_{-1},x_{0}$, and $x_{k+1}$ and%
\[
\varphi(T_{j},x_{j},u_{j})=\varphi(T_{j},x_{j},v_{j})\text{ for }%
j=0,\ldots,k-1.
\]
Hence (\ref{r1}) yields for $j=-2,\ldots,k$
\begin{align*}
d_{M}(\varphi(T_{j},x_{j},v_{j}),x_{j+1})  &  <\varepsilon_{M}(\varphi
(T_{j},x_{j},v_{j}))\\
&  \leq\varepsilon(v_{j}(T_{j}+\cdot),\varphi(T_{j},x_{j},v_{j}))=\varepsilon
(\Phi_{T_{j}}(v_{j},x_{j})).
\end{align*}
Concerning the component in $\mathcal{U}$ we have to show that%
\begin{equation}
d_{\mathcal{U}}(v_{j}(T_{j}+\cdot),v_{j+1})<\varepsilon(\Phi_{T_{j}}%
(v_{j},x_{j}))\text{ for }j=-2,\,-1,\ldots,k. \label{controls}%
\end{equation}
By choice of $T$ and $N$ all $w_{1},\,w_{2}\in\mathcal{U}$ satisfy%
\begin{align*}
&  d_{\mathcal{U}}(w_{1},w_{2})=\sum_{i=1}^{\infty}\frac{1}{2^{i}}%
\frac{\left\vert \int_{\mathbb{R}}\left\langle w_{1}(t)-w_{2}(t),y_{i}%
(t)\right\rangle \,dt\right\vert }{1+\left\vert \int_{\mathbb{R}}\left\langle
w_{1}(t)-w_{2}(t),y_{i}(t)\right\rangle \,dt\right\vert }\\
&  \leq\sum_{i=1}^{N}\frac{1}{2^{i}}\left\{  \left\vert \int_{\mathbb{R}%
\setminus\left[  -T,T\right]  }\left\langle w_{1}(t)-w_{2}(t),y_{i}%
(t)\right\rangle \,dt\right\vert \right.  +\left.  \left\vert \int_{-T}%
^{T}\left\langle w_{1}(t)-w_{2}(t),y_{i}(t)\right\rangle \,dt\right\vert
\right\}  +\widehat{\varepsilon}\\
&  <2\widehat{\varepsilon}+\max_{i=1,\ldots,N}\int_{-T}^{T}\left\vert
w_{1}(t)-w_{2}(t)\right\vert \,\left\vert y_{i}(t)\right\vert \,dt.
\end{align*}
For all considered pairs of control functions in (\ref{controls}) the
integrands vanish, hence the distances are bounded by $2\widehat{\varepsilon}%
$, and for $j=0,\ldots,k-1$ the definition of $\widehat{\varepsilon}$ shows
that%
\[
2\widehat{\varepsilon}=\min_{j=-2,\ldots,k}\varepsilon_{M}(\varphi(T_{j}%
,x_{j},v_{j}))\leq\min_{j=-2,\ldots,k}\varepsilon(\Phi_{T_{j}}(v_{j},x_{j})).
\]
Thus the claim in (\ref{controls}) holds.

(ii) Let $\mathcal{E}^{\ast}$ be a maximal invariant strongly chain transitive
set in $\mathcal{U}\times M$. Suppose that for $x,y\in M$ there are
$u,v\in\mathcal{U}$ with $(u,x),(v,y)\allowbreak\in E^{\ast}$, hence
$\varphi(t,x,u),\varphi(t,y,v)\allowbreak\in E^{\ast}$ defined by
(\ref{down_E}) for all $t\in\mathbb{R}$. Choose $\varepsilon_{M}\in
\mathcal{P}(M)$ and $T>0$ and define $\varepsilon\in\mathcal{P}(\mathcal{U}%
\times M)$ by $\varepsilon(u,x):=\varepsilon_{M}(x),u\in\mathcal{U},x\in M$. A
corresponding $(\varepsilon,T)$-chain in $\mathcal{U}\times M$ for the control
flow $\Phi$ yields $x_{j},u_{j},T_{j}$ such that the corresponding
trajectories satisfy%
\begin{align*}
d_{M}(\varphi(T_{j},x_{j},u_{j}),x_{j+1})  &  \leq d(\Phi_{T_{j}}(x_{j}%
,u_{j}),(u_{j+1},x_{j+1}))<\varepsilon(\Phi_{T_{j}}(x_{j},u_{j}))\\
&  =\varepsilon_{M}(\varphi(T_{j},x_{j},u_{j})).
\end{align*}
This shows that $E^{\ast}$ defined by (\ref{down_E}) is strongly chain controllable.

The proof of (i) and (ii) is concluded by the observation that $E^{\ast}$ is a
maximal invariant strongly chain controllable set if and only if
$\mathcal{E}^{\ast}$ is a maximal invariant strongly chain transitive set.
\end{proof}

\end{document}
